\section{Method}
\label{sec:method}

\begin{figure}[t]
  \centering
  \begin{tikzpicture}[
    font=\footnotesize,
    node distance=4.5mm and 2.7mm,
    box/.style={draw=inkTwo!60, rounded corners=2pt, fill=white, align=center,
                minimum height=10mm, inner sep=3pt, text width=#1},
    box/.default=17mm,
    hl/.style={box=#1, fill=softBlue, draw=sOne},
    hl/.default=17mm,
    arr/.style={-{Stealth[length=2mm]}, draw=inkTwo, line width=0.6pt},
    lab/.style={font=\scriptsize\itshape, text=inkTwo},
  ]
  % main row
  \node[box=12.5mm] (cam) {Webcam\\or clip};
  \node[box, right=of cam] (mp) {MediaPipe\\Holistic\\\scriptsize 543 points};
  \node[box, right=of mp] (sub) {Unified\\subset\\\scriptsize $T\times89\times3$};
  \node[hl=22mm, right=of sub] (pre) {\textbf{preprocess()}\\\scriptsize trim $\cdot$ mirror\\[-1pt]\scriptsize scale $\cdot$ resample};
  \node[box, right=of pre] (feat) {Features\\\scriptsize $64\times440$};
  \node[box, right=of feat] (model) {Conv-Transformer\\or BiGRU};
  \node[box=12mm, right=of model] (out) {Top-5\\glosses};
  \foreach \a/\b in {cam/mp, mp/sub, sub/pre, pre/feat, feat/model, model/out}
    \draw[arr] (\a) -- (\b);

  % training branch
  \node[box=17mm, below=8mm of pre, fill=soft] (aug) {augment()\\\scriptsize train only};
  \node[box=17mm, left=6mm of aug, fill=soft] (data) {ASL / LSFB\\datasets};
  \draw[arr] (data) -- (aug);
  \draw[arr] (aug) -- (pre);
  \node[box=17mm, below=8mm of model, fill=soft] (onnx) {ONNX FP32\\$\to$ INT8};
  \draw[arr] (model) -- (onnx);

  % groupings
  \begin{scope}[on background layer]
    \node[fit=(cam)(mp)(sub)(pre)(feat)(model)(out), inner sep=3.5mm, rounded corners=4pt,
          fill=sOne!4, draw=sOne!30] (rt) {};
  \end{scope}
  \node[lab, anchor=south west] at (rt.north west) {in the browser: MediaPipe Tasks JS $+$ preprocess.js $+$ ONNX Runtime Web};
  \node[lab, anchor=north west] at ($(data.south west)+(0,-1mm)$) {training (PyTorch)};
\end{tikzpicture}

  \caption{System overview. The same preprocessing function serves training (Python) and the
  browser runtime (JavaScript port, tested for parity). Only landmarks are processed; video never
  leaves the device.}
  \label{fig:pipeline}
\end{figure}

\Cref{fig:pipeline} summarises the system. A deterministic preprocessing function maps a raw
landmark sequence to a fixed-size feature matrix, which is augmented during training only and
classified by a sequence model.

\subsection{Preprocessing}
\label{sec:preprocessing}

Let $X\in\mathbb{R}^{T\times L\times3}$ be a raw sequence. Preprocessing applies the following
steps in order.

\paragraph{Temporal trimming.} Leading and trailing frames without any detected hand are removed;
a sequence in which no hand is ever detected is kept unchanged.

\paragraph{Canonical dominant hand.} One-handed signs are produced with the dominant hand, and the
handedness assigned by the landmark estimator depends on whether the image is mirrored. We count
the frames in which each hand slot is detected and, if the left slot is more frequent, mirror the
sequence: $x\mapsto 1-x$ for all points, combined with the permutation $\pi$ that swaps left and
right hands, shoulders, elbows, wrists and symmetric lip points. The resulting features are
invariant to a horizontal flip of the input whenever one hand dominates, a property verified by a
unit test.

\paragraph{Spatial normalisation.} Let $s^{l}_t$ and $s^{r}_t$ denote the shoulder positions in
frame $t$, and $\mathcal{T}_s$ the set of frames in which both are detected. All points are
centred and scaled as
\begin{equation}
  c = \frac{1}{|\mathcal{T}_s|}\sum_{t\in\mathcal{T}_s}\frac{s^{l}_t+s^{r}_t}{2}, \qquad
  w = \frac{1}{|\mathcal{T}_s|}\sum_{t\in\mathcal{T}_s}\lVert s^{l}_t-s^{r}_t\rVert_2, \qquad
  \tilde{X}_{t,j} = \frac{X_{t,j}-c}{w},
\end{equation}
applied to $(x,y)$, with $z$ divided by $w$. Both constants are computed over the whole clip, so
that the motion of the hands relative to the body is preserved. If the shoulders are missing or
degenerate ($w<0.02$), the centroid of all visible points and four times their mean standard
deviation are used instead.

\paragraph{Fixed length.} The clip is linearly resampled to $T'=64$ frames, which also normalises
signing speed and frame rate (30\,fps for ASL, 50\,fps for LSFB). When an interpolation involves a
missing neighbour, the value of the nearest frame is used. Zero-padding with a mask is evaluated as
an alternative in \cref{sec:ablations}.

\paragraph{Features.} Each frame is described by (i) the normalised $(x,y)$ coordinates of the 89
landmarks; (ii) the dominant hand expressed relative to its wrist and divided by the length of the
segment from the wrist to the base of the middle finger, which separates handshape from hand
position; and (iii) the first-order temporal differences of (i) and (ii). Depth is discarded by
default. This yields $F = 2\times(89+21)\times2 = 440$ features per frame. Remaining missing values
are set to zero, i.e.\ missing landmarks are encoded at the body centre.

\subsection{Data augmentation}
Augmentation is applied during training only, in two stages. The first stage acts on the raw
sequence: a time warp combining a global speed change ($\sigma\sim\mathcal{U}[0.8,1.2]$, the clip
being resampled to $T/\sigma$ frames) with a smooth monotonic warp of the time axis over four
segments with relative speeds in $[0.7,1.3]$; random frame dropping ($p=0.1$); random masking of
individual landmarks ($p=0.05$) and of a whole hand in a frame ($p=0.05$); and Gaussian noise
($\sigma=0.002$). The non-uniform component is essential: since every clip is resampled to $T'$
frames, a uniform speed change alone would be almost entirely undone by length normalisation.

The second stage is applied with probability $0.8$ \emph{after} shoulder normalisation, in the
body frame: a rotation of up to $\pm15^\circ$, a per-axis scaling in $[0.8,1.2]$ and a translation
of up to $0.1$ shoulder widths. Applied in image coordinates, before normalisation, translations
and isotropic scalings would be cancelled exactly, as our unit tests confirm. Random draws are
seeded by (seed, epoch, sample) for reproducibility.

\subsection{Models}

\paragraph{Input projection.} Both models first project each frame to $d=192$ dimensions with two
linear layers, layer normalisation and GELU activations.

\paragraph{BiGRU.} Two bidirectional GRU layers \citep{cho2014gru} with 192 hidden units per
direction, followed by masked mean and max pooling over time and a linear classifier
(\numGruParams{}\,M parameters).

\paragraph{Convolution--Transformer.} After a learned positional embedding, the network stacks two
stages, each made of three convolution blocks and one Transformer block
\citep{vaswani2017attention}. A convolution block is a pre-norm residual unit: pointwise expansion
with a gated linear unit \citep{dauphin2017glu}, depthwise temporal convolution with kernel size 17,
batch normalisation, SiLU and pointwise projection, as in the convolution module of the Conformer
\citep{gulati2020conformer}. Transformer blocks use four heads and an MLP ratio of two. The
sequence is mean-pooled under the mask before the classifier (\numTrParams{}\,M parameters).
Attention is implemented with plain matrix products, so that the model exports to ONNX without
custom operators. We refer to this model as the Conv-Transformer.

\subsection{Training}
Models are trained with AdamW \citep{loshchilov2019adamw} (learning rate $10^{-3}$, weight decay
$0.05$, batch size 128), a cosine schedule with three warm-up epochs, label smoothing of $0.1$
\citep{szegedy2016rethinking} and gradient clipping at norm 1, for at most 50 epochs. Early stopping
monitors the validation macro-F1 with a patience of 10 epochs, and the checkpoint with the best
validation macro-F1 is retained. Mixed precision is used on GPU. \Cref{app:hparams} lists all
hyper-parameters.

\paragraph{Transfer.} For cross-lingual transfer, all weights except the classifier are
initialised from the ASL model, and a new classifier is trained for the target vocabulary. Since
the input representation is identical across datasets, no adapter is required. We compare full
fine-tuning with a variant that freezes the encoder during the first five epochs. Two additional
controls test the robustness of the result: fine-tuning with a five times lower learning rate, and
a linear probe in which the encoder stays frozen throughout training.

\subsection{Browser runtime}
\label{sec:runtime}

The demo runs the MediaPipe Holistic Landmarker (Tasks API, WebAssembly with GPU delegate) on the
webcam stream, builds the 89-landmark frame and passes it to a segmenter that detects the start
and end of each sign. The segmenter is a state machine: a sign starts after three consecutive
frames with a visible hand and ends when the hands disappear for six frames, remain still (mean
hand displacement below $0.012$ shoulder widths per frame) for fourteen frames, or after 120
frames. The segment is then preprocessed by the JavaScript port and classified by the ONNX model
with ONNX Runtime Web \citep{onnxruntime} (WebAssembly backend).

Mismatches between training and deployment preprocessing are a common cause of failure in
deployed systems. We therefore generate test fixtures from the Python implementation (missing
hands, mirrored input, missing shoulders, single-frame and long clips, padding, depth features)
and verify that the JavaScript port reproduces the features within \numJsParity{}, i.e.\ at
single-precision accuracy. We also test the complete browser path: executed in headless Chromium
through the code of the demo, the LSFB test clips shipped with the demo receive the same top-1
prediction as with the PyTorch checkpoint (\numBrowserSame{} clips), with probabilities differing
by at most $\numBrowserDiff{}$.
